\scanpage{223}
Revenant au cas $A$ de car. $0$, que \unsure{puis|peut} on dire de $H^0_{DR}$ et $H_0^{\mathrm{Kos}}$ ? Pour $H_0$, on sait déjà que c'est $\simeq A$, est-il vrai que l'hom évident $A\longrightarrow H^0_{DR}$ est aussi un isom ? \deleted{\unsure{car}} i.e. que $\forall f\in\operatorname{Sym}^*(M)$ tel que $df=0$ soit \deleted{nulle} une constante, ? i.e. \inserted{si $h>0$,} $f\in\operatorname{Sym}^h(M)$, $h>0$ et $df=0\Longrightarrow f=0$, i.e.
\[
\operatorname{Sym}^hM\overset{d}{\unsure{-|\longrightarrow}}(\operatorname{Sym}^{h-1}M)\otimes\underset{\substack{\Vert\\\Lambda^1M}}{M}
\]
\margin{?}
\[
\overset{\unsure{\uparrow|\cap}}{(\operatorname{Sym}^{h-1}M)\otimes M}
\]
est injectif ? Or composant avec
\[
\operatorname{Sym}^{h-1}(M)\otimes\underset{\substack{\Vert\\\operatorname{Sym}^1M}}{M}\overset{\delta}{\unsure{-|\longrightarrow}}\operatorname{Sym}^hM
\]
on trouve la multiplication par $h$ (car $\theta=d\delta+\delta d$ est égal à $\delta d$ sur $\operatorname{Sym}^*M$ puisque $\delta$ y est nulle), donc inj, OK (Plus gén., ceci marche si $h_A$ inv. dans $A$.)

\uline{Dualité} Supposons $M$ et $M'$ deux $A$-modules accouplés vers $A$,
\[
M\otimes M'-A
\]
alors on en conclut des accouplements (degré par degré)
\[
\operatorname{Sym}^*M\times\Gamma^*M'-A
\]
\[
\Lambda^*M\times\Lambda^*M'-A
\]
qui est une dualité parfaite si $M$ est projectif t.f. et $M'$ s'identifie à son dual. On en conclut des accouplements
\[
\deleted{\operatorname{Sym}}S^*\Lambda^*(M)\times\Gamma^*\Lambda^*(M')\longrightarrow A
\]
Ceci posé, on vérifie que les op. $d_M$ et $\delta_{M'}$, $\delta_M$ et $d_{M'}$ sont \deleted{\unsure{alors}} en transposition pour cet accouplement. D'ailleurs, on vérifie aussi que l'op. de multiplication d'un côté est en dualité avec l'op. de comultiplication de l'autre. Donc dans le cas de la dualité parfaite, les opérations $d_M$ et $\delta_{M'}$, $\delta_M$ et $d_{M'}$, multiplication et comultiplication, se déterminent mutuellement.
\margin{222}
